% The receive framer. Two states, and every rejection has a counter.
\begin{tikzpicture}[font=\sffamily\scriptsize]
\node[initial] (i) at (0,0) {};
\node[state, text width=22mm] (hunt) at (2.4,0) {Hunting\\{\tiny discarding until a delimiter}};
\node[state, text width=22mm] (coll) at (7.8,0) {Collecting\\{\tiny appending to the frame buffer}};

\draw[trans] (i) -- (hunt);
\draw[trans] (hunt) -- node[above]{delimiter} (coll);

\draw[trans] (hunt.north west) to[out=125,in=55,looseness=1.8] (hunt.north east);
\node[anchor=south, align=center, font=\sffamily\tiny] at (2.4,1.75)
  {any other byte:\\discard it};

\draw[trans] (coll.north west) to[out=125,in=55,looseness=1.8] (coll.north east);
\node[anchor=south, align=center, font=\sffamily\tiny] at (7.8,1.75)
  {any other byte:\\append it};

\draw[trans] (coll.north east) to[out=35,in=-35,looseness=3.5] (coll.south east);
\node[anchor=west, align=left, font=\sffamily\tiny, text width=28mm] at (10.1,0)
  {delimiter: complete the frame, then start the next one};

\draw[trans] (coll.south) to[out=250,in=290,looseness=1.0] (hunt.south);
\node[anchor=north, font=\sffamily\tiny] at (5.1,-1.45) {buffer full: oversize++};

\node[umlnote, text width=44mm, anchor=north west] at (-0.4,-2.2)
  {Completing a frame runs three checks in order, cheapest first, and each one
   has its own counter: bad stuffing, too short, bad checksum. A link that is
   failing then tells you how it is failing and not only that it is.};
\node[umlnote, text width=44mm, anchor=north west] at (5.4,-2.2)
  {There is no state for a partly received length and no guessing where the
   next frame might begin. That is the whole of what the delimiter bought, and
   it is why this machine has two states rather than five.};
\end{tikzpicture}
